{\scriptsize\setlength{\tabcolsep}{4pt}
\begin{longtable}{>{\raggedright\arraybackslash}p{0.07\textwidth}>{\raggedright\arraybackslash}p{0.41\textwidth}>{\raggedright\arraybackslash}p{0.17\textwidth}>{\raggedright\arraybackslash}p{0.25\textwidth}}
\caption{Results that the main text (by section) and the supplement quote from runs or computations that no other table or figure shows, and the inputs that the earlier study set without a primary source (Table~2 of the main text, $^{a}$), each with the runs or the code it comes from. Every value here is read from its file by the numbers check's own entry for it.}\label{tab:s-quoted}\\
\toprule
Where & Result & Value & From \\
\midrule
\endfirsthead
\multicolumn{4}{l}{\textit{Continued}} \\
\toprule
Where & Result & Value & From \\
\midrule
\endhead
3.1 & With body energies up to 33.2 kJ and no delivered heat, over the 5,248 runs at every contact, tab conductance, release time and connection and under the 200 kinetic draws: the highest peak of the resolved neighbour's hottest point, its margin to 255.7 \textdegree C, and the lumped neighbour's highest peak & 226.7 \textdegree C; 29.0 K; 180.6 \textdegree C & step 5's runs; \nolinkurl{results/v12_numbers.json}, all\_runs \\
3.1 & The eight tests with the largest fitted triggers held to 33.2 kJ, no delivered heat: the hottest point's peak at the high contact (two connections) and the lumped neighbour's highest peak & 226.4 and 226.7 \textdegree C; 177.4 \textdegree C & step 12's runs; \nolinkurl{results/stage12_paper_numbers.json}, part\_a.high \\
3.1 & With the fitted triggers and no delivered heat, the latest time by which conduction carries the hottest point over 255.7 \textdegree C in the eight tests with the largest triggers & 21.7 s & step 12's runs; \nolinkurl{results/stage12_paper_numbers.json}, part\_a.central \\
3.1 & With the fitted triggers and no delivered heat at the low contact, the highest peak of the hottest point over the sixteen tests & 162.6 \textdegree C & step 7's runs; \nolinkurl{results/stage7_crossings.csv} \\
3.3 & In the three-dimensional threshold searches of PT-Exp4, PP-Exp2 and SP-Exp3: the lumped neighbour's runaway times, its reaction heat by 900 s, and the resolved neighbour's reaction heat by its own runaway & 106 to 128 s; 22.4 to 22.6 kJ; 4.1 to 4.8 kJ & step 15's runs; \nolinkurl{results/stage15_paper_numbers.json}, B1 \\
3.3 & ST-Exp3 with 10.80 kJ over 21.4 s: the resolved neighbour's runaway & 61.5 s & step 15's run; \nolinkurl{results/stage15_paper_numbers.json}, B1 \\
3.4 & PT-Exp1 with a 360 s pulse at the upper end of the local threshold's bracket: the hottest point's crossing, on the cell's axis, and the mean then & 316.5 s; 240.1 \textdegree C & step 13's run; \nolinkurl{results/stage13_skin.csv} \\
3.4 & The lumped neighbour's ignition temperature with the trigger at 25 \textdegree C, and the share of the low-temperature cathode reaction in the reaction heat at the ignition temperature & 230.4 \textdegree C; 73.3 \% & step 13's computation; \nolinkurl{results/stage13_tign.csv} \\
3.4 & With no delivered heat in PT-Exp4, PP-Exp3 and ST-Exp2 at the central contact, the heat the lumped neighbour has taken in through the contact and wedge by 250 s beyond the resolved one's & 2.1 to 2.3 kJ & step 15's accounts of archived runs; \nolinkurl{results/stage15_accounts.json} \\
3.4 & PT-Exp4 with 5.04 kJ over 21.4 s: the lumped neighbour's runaway under the kinetic draws ranked 5 and 50 & 191.0 and 196.5 s & step 15's runs; \nolinkurl{results/stage15_paper_numbers.json}, B4 \\
3.4 & PT-Exp4 with 5.04 kJ over 21.4 s and every rate law shifted by 20 K: the resolved neighbour's runaway, and the far face's reading where its jump begins & 270.3 s; 240.4 \textdegree C & step 15's run; \nolinkurl{results/stage15_paper_numbers.json}, B4 \\
3.5 & PT-Exp4, the 5.04 kJ pulse followed by 10 W: the resolved neighbour's runaway & 369.3 s & step 15's run; \nolinkurl{results/stage15_paper_numbers.json}, B3b \\
3.5 & PT-Exp4: the root mean square of the far face's difference from the record over 0 to 271.5 s with 4.66 and 5.04 kJ over 21.4 s, and over 0 to 181.9 s with 7.42 kJ & 16.3 and 16.8 K; 51.4 K & steps 10 and 15; \nolinkurl{results/stage10_paper_numbers.json}, \nolinkurl{results/stage15_paper_numbers.json} \\
3.5 & PT-Exp4 with 7.42 kJ over 21.4 s, read post hoc: the two side probes' readings where the far face's jump begins & 235.5 and 235.8 \textdegree C & step 15's run; its series, read by its independent read (\nolinkurl{results/stage15_independent_read.md}) \\
3.7 & Wall time of a run on one two-core machine: the reactive cross-section, and the three-dimensional model in the same searches & 17 to 71 s; 246 to 512 s & \nolinkurl{results/stage15_xs/*.json} and \nolinkurl{results/stage15_runs/*.json}, wall\_s \\
S2 & Step 3, the earlier geometry's case (33.2 kJ over 30 s, a 500 mW\,K$^{-1}$ busbar): the hot spot's onset over the 200 kinetic draws, under the four draws step 11 took, under the posterior mean, and its standard deviation over the kinetic and over the 30 thermal draws & 6.01 to 17.13 s; 6.25, 10.23, 16.04, 17.02 s; 11.48 s; 2.83 and 0.076 s & step 3's runs; \nolinkurl{results/stage3a_onsets.csv}, stage3a\_metrics.json, stage11\_draws.csv, stage3bc\_metrics.json \\
Table 2 & Tab conductance per tab, set by the earlier study without a primary source and held here as an assumption & 6.8 mW\,K$^{-1}$ (5 and 10 as brackets) & \nolinkurl{src/v12_model.py} line 37 (G\_TAB = 0.005, 0.0068, 0.010 \nolinkurl{W/K}) \\
Table 2 & Can conductivity, set by the earlier study without a primary source and held here as an assumption & 15 W\,m$^{-1}$\,K$^{-1}$ & \nolinkurl{external/m50-two-cell-pilot/src/lumped_two_cell.py} line 57 (K\_CAN = 15.0) \\
Table 2 & Can mass and specific heat, set by the earlier study without a primary source and held here as an assumption & 9 g; 500 J\,kg$^{-1}$\,K$^{-1}$ & \nolinkurl{external/m50-two-cell-pilot/src/lumped_two_cell.py} lines 54 to 55 (M\_CAN, CP\_CAN) \\
Table 2 & Air in the wedge, set by the earlier study without a primary source and held here as an assumption & 0.03 W\,m$^{-1}$\,K$^{-1}$ & \nolinkurl{external/m50-two-cell-pilot/src/lumped_two_cell.py} line 69 (K\_AIR = 0.03) \\
Table 2 & Contact sector, either side of the line of contact, set by the earlier study without a primary source and held here as an assumption & 45\textdegree & \nolinkurl{src/v12_mesh.py} line 22 (SECTOR = pi / 4, "half-angle" \\
Table 2 & Room convection, set by the earlier study without a primary source and held here as an assumption & 10 W\,m$^{-2}$\,K$^{-1}$ & \nolinkurl{external/m50-two-cell-pilot/src/lumped_two_cell.py} lines 63 and 74 (H\_CONV, A\_CONV) \\
Table 2 & Cell mass, set by the earlier study without a primary source and held here as an assumption & 65 g & \nolinkurl{external/m50-two-cell-pilot/src/lumped_two_cell.py} line 47 (M\_CELL = 0.065) \\
\bottomrule
\end{longtable}
}
